\chapter{Recolha, preparação e tratamento dos datasets}
\label{chap:data}

Este capítulo descreve a componente atribuída ao módulo de preparação de dados do projeto
Basketball AI Coach: recolha, preparação e tratamento dos conjuntos de
dados (\emph{datasets}). Esta componente integra AAUTIA e fornece uma base
comum aos módulos de aprendizagem automática e de planeamento de PLANTIA.

A execução documentada em \path{docs/DATA_PREPARATION.md} e
\path{docs/DATA_QUALITY_REPORT.md} do projeto \texttt{challenge-two}
preparou 8\,283 jogos da fase regular da NBA, de 2019--20 a 2025--26,
com origem principal na coleção \texttt{kaggle\_eoinamoore\_historical\_nba}.
O relatório regista 16\,566 observações de equipa por jogo, 215\,473
registos de jogador por jogo e 4\,607\,885 eventos.

As secções seguintes descrevem a preparação realizada e os critérios
adotados, distinguindo as extensões ainda condicionadas à validação das fontes.
Na implementação documentada, os eventos estão preparados, mas os
alinhamentos em campo não foram reconstruídos nem certificados. O
módulo de avaliação de jogadores e alinhamentos utiliza os cinco iniciais observados; os minutos da fonte e os
fusos horários dos relatórios de lesões não foram validados para os usos
preditivos previstos. Os capítulos dos módulos de aprendizagem automática e de planeamento especificam os
subconjuntos e as transformações efetivamente utilizados.

\section{Responsabilidades e articulação com os restantes elementos}

O módulo de preparação de dados foi responsável pela seleção e recolha dos dados, limpeza,
análise exploratória, construção dos conjuntos de trabalho e documentação
das transformações. Couberam-lhe também os testes de qualidade e as
experiências de preparação, nomeadamente a comparação da cobertura e das
consequências de diferentes regras de tratamento de valores em falta.
Os requisitos dos dados, as interfaces e os critérios de avaliação foram
definidos em conjunto com os responsáveis pelos módulos que os utilizam.

A preparação foi orientada pelas seguintes necessidades:
\begin{itemize}
\item \textbf{Previsão de resultados:} dados por jogo,
indicadores históricos das equipas e resultado usado como variável-alvo.
\item \textbf{Avaliação de jogadores e alinhamentos:}
estatísticas individuais, minutos de utilização e, quando disponíveis,
intervalos de jogo com identificação dos jogadores em campo.
\item \textbf{Cinco inicial e distribuição de minutos:}
plantéis, posições e informação disponível sobre elegibilidade e utilização.
\item \textbf{Planeamento das substituições:} sequência
de eventos e substituições que permita construir estados e observar
transições, quando a fonte ofereça essa granularidade.
\item \textbf{Objetivos ao longo de vários jogos e
integração:} calendário, resultados históricos e identificadores comuns
para relacionar as entradas e saídas dos diferentes módulos.
\end{itemize}

O treino dos modelos, a estimação das transições e a definição dos
problemas de planeamento ficaram a cargo dos respetivos elementos.
O módulo de preparação de dados disponibilizou os dados e as regras de preparação necessários
a essas tarefas, distinguindo os registos observados das hipóteses
introduzidas pelos módulos.

\section{Seleção das fontes e do âmbito}

A seleção começou pela comparação de fontes que disponibilizavam
resultados, estatísticas de equipas e jogadores e, quando possível,
registos de eventos \emph{play-by-play}. Para cada fonte foram registados
a cobertura temporal, as competições disponíveis, a granularidade,
a documentação dos campos, as condições de acesso e reutilização e a
possibilidade de reproduzir a recolha. Foi privilegiada uma competição
com várias temporadas consecutivas e identificadores consistentes.

A escolha atendeu às necessidades de cada tarefa. A previsão
pré-jogo requer resultados e estatísticas históricas; a avaliação de
alinhamentos exige identificar os cinco jogadores em campo e os intervalos
da sua utilização. Dados agregados por partida podem sustentar o primeiro
objetivo sem permitir o segundo. Antes da seleção definitiva, uma amostra
foi inspecionada para confirmar a existência e a qualidade dos campos
necessários. Como os alinhamentos não foram reconstruídos com confiança,
a avaliação utilizou os cinco iniciais observados, mantendo explícita
essa limitação.

Foram definidos os critérios de inclusão de jogos, temporadas e fases
da competição. Partidas canceladas, registos incompletos e alterações
de formato foram tratados segundo regras documentadas. A quantidade
final de jogos e jogadores, o período abrangido e os motivos de exclusão
foram apresentados após a recolha.

\section{Recolha e organização dos dados}

Os ficheiros originais foram preservados, acompanhados pela identificação
da fonte, data de recolha e versão do procedimento utilizado. A organização
separou dados originais, dados limpos e conjuntos preparados para
modelação, permitindo repetir as transformações sem modificar a origem.

A estrutura relacionou jogos, equipas e jogadores através de
identificadores estáveis. Os jogos incluíram data, competição, temporada,
equipas e resultado; as estatísticas individuais foram associadas ao
jogador, à equipa e à partida. Quando disponíveis, os eventos incluíram
o período, o instante, o tipo de ocorrência e os intervenientes.
Um dicionário de dados descreveu o significado, a unidade e os valores
admissíveis de cada campo.

A integração entre fontes exigiu verificar correspondências de nomes,
identificadores, datas e equipas. Transferências de jogadores foram
representadas com referência temporal, evitando atribuir a uma partida
informação sobre um plantel posterior. A ordenação dos eventos respeitou
o período e o sentido do relógio de jogo, incluindo prolongamentos.

\section{Limpeza e análise exploratória}

A limpeza identificou duplicados, campos em falta, tipos incorretos e
valores inconsistentes. Foram verificadas, por exemplo, a unicidade dos
jogos, a relação entre lançamentos convertidos e tentados e a coerência
dos resultados. Um valor ausente foi distinguido de um zero observado,
como no caso de estatísticas de jogadores que não participaram.

O tratamento de omissões dependeu do campo e da tarefa. A ausência do
resultado impede a utilização de uma partida como exemplo supervisionado;
a falta de uma variável explicativa podia justificar imputação, um
indicador de ausência ou exclusão. As decisões foram documentadas e a
sua incidência quantificada. Valores extremos foram inspecionados antes
de qualquer decisão de remoção, pois podem corresponder a desempenhos reais ou a
jogos com prolongamento.

A análise exploratória caracterizou a cobertura por temporada, os dados
em falta, a distribuição dos resultados e as diferenças entre equipas.
Foram também examinadas mudanças de distribuição ao longo do tempo.
A exploração usada para escolher variáveis e transformações ficou
limitada aos dados de desenvolvimento, preservando o conjunto de teste
para a avaliação final.

\section{Construção do conjunto para previsão de resultados}

Para a previsão pré-jogo, cada observação representou uma partida,
com características das duas equipas e uma variável-alvo binária:
vitória da equipa da casa ou da equipa visitante. Jogos em campo neutro
tiveram uma indicação própria. O resultado final serviu para construir
o alvo, mas não integrou as variáveis explicativas da mesma partida.
Identificadores foram conservados para rastreabilidade e agrupamento,
sem serem automaticamente utilizados como entradas do modelo.

As características candidatas incluíram médias móveis de pontos,
eficiência ofensiva e defensiva, aproveitamento dos lançamentos,
ressaltos, assistências, perdas de bola e dias de descanso. Podiam ser
representadas separadamente por equipa ou por diferenças entre as duas.
As janelas históricas foram definidas durante a validação, e o início
de temporada ou a falta de histórico tiveram um tratamento explícito.

Todas as agregações foram calculadas apenas com jogos anteriores ao
instante da previsão. Estatísticas da própria partida, classificações
finais da temporada e informação posterior sobre lesões ou alinhamentos
não podiam alimentar uma previsão pré-jogo. A utilização de informação
sobre jogadores disponíveis dependeu de ser possível confirmar que
estava acessível nesse instante.

As variáveis categóricas foram codificadas e as numéricas podiam ser
normalizadas quando o algoritmo o justificava. Imputação, normalização,
seleção de características e outras transformações aprendidas foram
ajustadas exclusivamente nos dados de treino de cada divisão. O mesmo
procedimento foi depois aplicado à validação e ao teste, sem recalcular
parâmetros nesses conjuntos.

\section{Construção do conjunto para avaliação de jogadores}

O conjunto individual teve uma observação por jogador e partida, associada
à equipa que representava nessa data. Incluiu, conforme a cobertura da
fonte, minutos, pontos, lançamentos, ressaltos, assistências, perdas de
bola e faltas. A participação efetiva foi distinguida da presença no
plantel, para não interpretar a ausência de utilização como um desempenho
nulo observado em campo.

Foram preparados indicadores históricos de utilização e desempenho,
calculados apenas a partir de partidas anteriores ao instante de previsão.
A comparação de jogadores podia recorrer a estatísticas por minuto ou
por posse, quando existiam dados suficientes para calcular esses valores.
Denominadores nulos e amostras reduzidas tiveram tratamento explícito,
conservando-se o número de jogos e os minutos que sustentam cada indicador.

A variável-alvo de desempenho foi acordada com o módulo de avaliação de jogadores e alinhamentos, identificando
a unidade de observação e o horizonte a prever. As estatísticas da partida
avaliada podiam servir como alvo ou como medida de avaliação, mas não
como informação de entrada para prever essa mesma partida. O protocolo de preparação
dos intervalos de alinhamentos é descrito na secção dedicada ao planeamento.

\section{Divisão temporal e prevenção de fuga de informação}

O conjunto foi dividido cronologicamente em treino, validação e teste.
Os jogos mais antigos serviram para aprender os modelos, um período
posterior para selecionar configurações e o período mais recente para
a avaliação final. As datas de corte dependeram da cobertura disponível
e foram registadas antes da comparação final, juntamente com o número
de observações em cada divisão.

O protocolo de desenvolvimento contemplou a possibilidade de validação
com janela de treino crescente, com aprendizagem sobre o passado e
avaliação num bloco posterior. Não foi usada uma divisão aleatória que misturasse
jogos futuros e passados. Quando existiam várias observações da mesma
partida, todas permaneceram na mesma divisão.

O protocolo distinguiu a atualização de indicadores com resultados de
jogos entretanto concluídos do reajuste do modelo. Numa avaliação
sequencial, esses resultados podiam atualizar o histórico para o jogo
seguinte, desde que isso correspondesse ao funcionamento previsto e não
fosse usado para escolher configurações a partir do desempenho no teste.

\section{Dados para alinhamentos e planeamento}

A reconstrução de intervalos com um alinhamento constante não foi
realizada na execução documentada. O protocolo definido para essa extensão
previa registar jogadores, duração, adversário e diferencial de pontos,
conhecer o cinco inicial e processar substituições pela ordem correta.
Previa também sinalizar intervalos ambíguos ou incompatíveis com cinco
jogadores por equipa e tratá-los segundo critérios explícitos.

Para essa extensão, ficaram definidos os seguintes requisitos: considerar
o tempo ou o número de posses observadas na avaliação dos alinhamentos,
manter os intervalos da mesma partida juntos na divisão temporal e
estimar desempenhos e transições apenas com os dados de treino.
Na execução realizada, a avaliação dos alinhamentos utilizou os cinco
iniciais observados, sem os equiparar a intervalos reconstruídos.

Para o planeamento foram ainda necessários calendário, plantel elegível
e restrições de utilização. Disponibilidade observada e hipóteses de
simulação foram distinguidas. A fadiga, se modelada a partir de minutos
e descanso, foi apresentada como uma aproximação definida no projeto,
e não como uma medição direta da condição física.

\section{Rastreabilidade e resultados da preparação}

O processo produziu um conjunto para previsão pré-jogo, um conjunto para
avaliação individual e dados de eventos. A construção de um conjunto de
intervalos de alinhamentos permaneceu pendente de reconstrução e validação. Foram ainda
organizados os dados de calendário e plantel necessários à integração. Cada versão incluiu o dicionário de dados, as regras de
limpeza, a definição das características, os períodos de treino,
validação e teste e um resumo dos registos excluídos.

As interfaces de entrega especificaram o formato dos ficheiros, as chaves
de ligação, os tipos e unidades dos campos, a representação de valores
ausentes e o instante de referência da informação. Os conjuntos partilharam
os identificadores de jogo, equipa e jogador e um protocolo temporal
comum, permitindo cruzar previsões e decisões sem misturar períodos de
treino e avaliação. Alterações de estrutura ou de regras de preparação
foram registadas numa nova versão e comunicadas aos restantes elementos.

O relatório da preparação apresentou a cobertura por temporada, os
números de jogos e jogadores, a proporção de valores em falta e os efeitos
das exclusões e transformações. A entrega foi acompanhada pelos
procedimentos de recolha e preparação, pelas configurações utilizadas e
pelos resultados dos testes de qualidade, permitindo aos restantes
elementos reproduzir os conjuntos recebidos.

Os procedimentos e configurações foram conservados para permitir repetir
a preparação. Antes da modelação, foram verificadas a ausência de
sobreposição entre divisões, a coerência dos identificadores e a
disponibilidade temporal das variáveis. Esta documentação permitiu
relacionar os resultados experimentais com os dados e as hipóteses que
os produziram.
